\scanpage{628}
\mspage{16}
\[
\left\{\begin{aligned}
\text{a)}\quad &\text{Factorisation pour } C'\longleftrightarrow TF'\text{ et }TC'\longleftrightarrow F',\\
&\text{i.e. les relations } C' \mathrel{\longleftarrow\!\mid\!\longrightarrow} TF'\text{ et }TC' \mathrel{\longleftarrow\!\mid\!\longrightarrow} F'.\\
\text{b)}\quad &TC'\subset W'.
\end{aligned}\right.\tag*{(2.1.4)}
\]

ou encore, au lieu de b), la condition équivalente
(en présence de a)) :

\[
\left\{\begin{aligned}
\text{b')}\quad &\text{Toute flèche }f'\text{ dans }M'\text{ (ou seulement}\\
&f'\in TC')\text{ se factorise en }f'=p'i',\text{ avec}\\
&i'\in W',\ p'\in F'.
\end{aligned}\right.\tag*{(2.1.5)}
\]

$2^\circ)$ Supposons \inserted{(2.1.4.} a)) satisfaite \inserted{(factorisation)}. Alors, pour que
$(W',C',F')$ soit un triple de Quillen clos,
il \uline{suffit} que l'on ait

\[
\text{b'')}\quad\left|\underline{u^*u_!(TC)\subset TC}\right|,\tag*{(2.1.6)}
\]

et que de plus on soit dans l'un des
trois cas suivants

A) $u^*$ admet un adjoint à droite $u_*$

B) $u^*$ commute aux \deleted{\unsure{pet}}\unsure{petites} $\varinjlim$, $(W,C,F)$
est spécial, et $M'$ est spécial i.e. satisfait à la
condition (2.6 a).

C) $u^*$ est pleinement fidèle, et $u^*u_!(F)\subset F$.

Notons tout de suite le

\uline{Corollaire} 2.2. Supposons que $(W,C,F)$ soit un
\deleted{\unsure{triple}} de Quillen spécial, que $M'$ soit spécial,
\deleted{et qu'on ait} que $u^*$ commute aux petites $\varinjlim$,
enfin que l'on ait (2.1.6) i.e. $u^*u_!(TC)\subset TC$.
Alors $(W',C',F')$ est un couple de Quillen
spécial.
% archive sheet number: 627
